\documentclass[10pt,letterpaper]{article}
\usepackage{cornell-notes}

\title{Clause 07.06.05: Multiplicative Operators}
\author{}
\date{}
\setDocTitle{Clause 07.06.05: Multiplicative Operators}
\setDocSubtitle{C++ 2024}
\setDocOwner{C++ 2024 Cornell Notes}
\setDocDate{\today}
\setCornellCollection{C++ 2024 Cornell Notes}
\setCornellUnitType{Clause}
\setCornellUnitNumber{07.06.05}
\setCornellUnitTitle{Multiplicative Operators}
\hypersetup{pdftitle={Clause 07.06.05: Multiplicative Operators},pdfsubject={Cornell notes},pdfauthor={}}

\begin{document}
\maketitle
\begin{CornellOverviewBox}{Overview}
\textbf{Classification:} Core language - Normative\\[2pt]
\textbf{Learning objective:} Understand the rules, terminology, and semantic consequences associated with Multiplicative operators.
\end{CornellOverviewBox}\begin{CornellSummaryBox}{Summary}
{\sffamily\bfseries\color{Accent} Summary}\par\vspace{3pt}Section 7.6.5 focuses on Multiplicative operators. Use the listed grammar productions together with the semantic constraints and disambiguation rules referenced by the section. When applying it, preserve the distinction between mandatory requirements, recommendations, permissions, and behavior deliberately left undefined, unspecified, or implementation-defined.
\end{CornellSummaryBox}

\end{document}
